Lately, my research has been focused on the question: what could be the requirements of any brain-like reasoning system? This question can be analyzed from the perspectives of philosophy, biology, genetics, mathematics and many others. For a lot of reasons that I will not expand here, I think that any living being, in order to act in the world and survive, can only rely on two things: million years of experience captured in DNA and sensory information coming in its way to quickly adapt to changing and threatening environments. The scope of this paper is narrowed down to the latter case because I argue that the process is likely similar for both but at different timescales and hypothetically justified by evolution theory.

Consequently, I assume that, a living being, also called agent or creature in the rest of this paper, only has its sensory evidences that it can act on in order to evolve and thrive in its environment. This concept of evidence, which somehow relates to proofs of a given statement and more generally to the concept of explanation, helped me make the connection with proof theory and type theory. Indeed, in those theories, any term belonging to a type is an evidence or a proof of a proposition. In type theory, propositions are usually encoded using a language derived from the language of pure \lambdacalc. Nowadays, there are many new such languages that assist mathematicians to build and formalize mathematical statements and proofs by making machines handle the burden of following mathematical reasoning meticulously. To name only a few of those languages, there is the \textit{Calculus of Construction} used by \textit{Coq}, that superseeded the simply-typed \lambdacalc, and there is also the language of Dependent types used by Lean and other proof assistants. From there, one question that begs attention is: could symbolic reasoning be anything else than mathematical reasoning or an approximation of it?

So far we know that the only thing that an agent has at hand to build its formal theory of explanation of the world is its raw sensory information. It is the only evidence available to it. Therefore, I hypothesize that any symbol entering the system, be it an analogic value or any other kind of digital symbol, must be the only raw material or evidence the agent has to build its theory of the world. Moreover, the connection with proof and type theories make the \lambdacalc language suitable to encode the statements the agent needs to build the theory. Some of the most advanced research on type theory is Dependent type theory and Homotopy type theory which paved the way to many proof assistants as already mentioned. However, it suffices to look at the Dependent type theory, and more particularly at the interplay between types and the body of the statement to see the potential connection between them. This delimitation line blurs quite a bit when we consider that terms can depend on types and types can depend on terms. Does it make sense to have a typing statement in addition to the statement already written by the \lambdaexpr?
I argue that this distinction was needed for two reasons:
\begin{itemize}
    \item We had the theory of \lambdacalc already and we did not want to go over it again with modified rules so we preferred creating a complement to the existing theory that ended up being two different theories.
    \item We, intelligent beings, need to bridge the gap between the proofs and the statements with some kind of subjective interpretation of the expression hidding the complexity of the algorithm we have in mind. In a nutshell, as mathematicians, we must hypothesize interesting statements, encode them and then build a transformation function that connects previously established proofs into the proofs of the new statements. However, this transformation can be more or less complex and subjective and this complexity and subjectivity gap must be filled by human intervention for lack of better intelligent system or algorithm being able to fulfill this task. For instance, $0$ has been chosen to be the $0$ of natural number. One might ask why $0$ and not the word \textit{dog} instead? Then $succ(dog)$ would represent what we usually call $1$. Where does this $0$ come from? From an intelligent system perspective, we cannot assume that it has the same axioms as the majority of mathematicians as long as we do not communicate it. The zero of natural number could theoretically be any symbol or combination of symbols coming from its sensory inputs. Therefore, I hypothesize that reasoning in an agent cannot rely on any kind of interpretation since it is not omniscient and therefore any incoming symbol or combination has as much chance as others to be considered as the zero of natural numbers until the assumption proves to be useful for explaining the observations.
\end{itemize} 
The last development then brings a question to the front: what is in between terms and types? It certainly is an interpretation that hides something deeper, which is that there exists a function that transforms the term into its type and that we do not know yet.


My proposal answer to that question is it must be something very simple and intrinsic to the model. Can a term becomes a type of its own and when does that happen? Is there a good reason to treat terms and types differently?
More questions begging a proposition of answer leading to the introduction of this new framework called \lambdatypes.